% 鞍点几何（Lagrangian）：用等高线 + 方向箭头示意鞍点结构
\documentclass[tikz,border=8pt]{standalone}
\usepackage{ctex}
\usepackage{amsmath}
\usepackage{pgfplots}
\pgfplotsset{compat=1.18}
\usetikzlibrary{calc,arrows.meta,3d}

\begin{document}
\begin{tikzpicture}

%% === 鞍点曲面：用 pgfplots 3D 表面图 ===
\begin{axis}[
  width=10cm, height=8cm,
  view={35}{25},
  xlabel={$x$（原始变量）},
  ylabel={$\lambda$（对偶变量）},
  zlabel={$L(x,\lambda)$},
  xlabel style={font=\small},
  ylabel style={font=\small},
  zlabel style={font=\small, rotate=-90},
  xtick={-2,-1,0,1,2},
  ytick={-2,-1,0,1,2},
  xmin=-2.5, xmax=2.5,
  ymin=-2.5, ymax=2.5,
  zmin=-3, zmax=5,
  colormap={mymap}{
    color(0)=(blue!30!white)
    color(0.5)=(white)
    color(1)=(red!30!white)
  },
  mesh/rows=25,
  mesh/cols=25,
  axis lines=box,
  tick label style={font=\scriptsize},
  title style={font=\small\bfseries},
  title={Lagrangian 鞍点曲面 $L(x,\lambda) = x^2 + \lambda(x-1)$\\极小化对 $x$，极大化对 $\lambda$},
]
%% L(x,lambda) = x^2 + lambda*(x-1)
%% 鞍点：对x最小: 2x+lambda=0 => x=-lambda/2，代入得最小值 -lambda^2/4 - lambda
%% 对lambda最大: g(lambda) = -lambda^2/4 - lambda，g'=0 => lambda*=-2，x*=1
%% 鞍点在 (x*,lambda*) = (1, -2)，L=1+(-2)(0)=1

\addplot3[
  surf,
  shader=faceted interp,
  opacity=0.75,
  domain=-2.2:2.2,
  domain y=-2.5:2.5,
  samples=20,
] {x^2 + y*(x-1)};

%% 鞍点位置标注（x=1, lambda=-2）
\addplot3[
  only marks, mark=*, mark size=3pt, color=red!80!black
] coordinates {(1, -2, 1)};

\end{axis}

%% === 鞍点说明文字（在图外添加）===
\node[draw=red!50, fill=red!5, rounded corners, font=\small,
      align=left, inner sep=5pt] at (9.5, 4.5)
  {\textbf{鞍点 $(x^*, \lambda^*)$}\\
   对 $x$ 取极\textbf{小}（下凸）\\
   对 $\lambda$ 取极\textbf{大}（上凸）\\[4pt]
   $\min_x \max_\lambda L = \max_\lambda \min_x L$\\
   （强对偶成立时）};

%% === 方向说明 ===
\node[font=\scriptsize, blue!70!black] at (0.5, 1.0)
  {$x$ 方向：凹形曲线（极小）};
\node[font=\scriptsize, red!70!black] at (0.5, 0.5)
  {$\lambda$ 方向：凸形曲线（极大）};

\end{tikzpicture}
\end{document}
